\section{Conclusion}
We introduced GaitEval, a six-dataset benchmark for grounding dataset-specified gait-label predictions in pose-derived motion evidence. Controlled inputs and same-feature linear references distinguish limited signal from failed evidence use. The five-family and same-model results retain the accepted finding of limited zero-shot agreement in many tested settings, with dataset--representation exceptions under the shared six-dataset protocol. Task-specific adaptation recovers selected in-distribution mappings; the retained ATAXIA/PATHOGAIT adaptation results do not establish clinical or population-level superiority. Reliable transfer is not observed in the tested bidirectional ATAXIA--PATHOGAIT probe. Overall, GaitEval shows that LLM grounding in pose-derived gait evidence is heterogeneous and representation-dependent: adaptation can recover dataset-specific mappings in distribution, but reliability beyond the tested settings remains unresolved.
